\documentclass[11pt,a4paper]{article}

\usepackage{amsmath}
\usepackage{amssymb}
\usepackage{amsthm}
\usepackage[margin=3cm]{geometry}
\usepackage{hyperref}

\theoremstyle{plain}
\newtheorem{theorem}{Theorem}
\newtheorem{proposition}[theorem]{Proposition}
\newtheorem{lemma}[theorem]{Lemma}
\newtheorem{corollary}[theorem]{Corollary}
\theoremstyle{definition}
\newtheorem{definition}[theorem]{Definition}
\theoremstyle{remark}
\newtheorem{remark}[theorem]{Remark}

\newcommand{\R}{\mathbb{R}}
\newcommand{\N}{\mathbb{N}}
\newcommand{\Q}{\mathbb{Q}}
\newcommand{\dom}{\operatorname{dom}}
\newcommand{\sem}[1]{[\![#1]\!]}
\newcommand{\fine}{\mathrm{fine}}
\newcommand{\can}{\mathrm{canon}}
\newcommand{\Sig}{\mathcal{S}}
\newcommand{\Ax}{\mathcal{A}}
% \bigsqcap is not standard (LaTeX ships \bigsqcup only); make one
% from \sqcap as a large operator, so no extra package is required.
\newcommand{\bigsqcap}{\mathop{\sqcap}\limits}

\title{The canonical signal space of an STL formula}
\author{TABEX}
\date{}

\begin{document}
\maketitle

\begin{abstract}
We show that two equivalent formulas of a bounded, discrete fragment of Signal
Temporal Logic have the same \emph{canonical signal space}, and conversely. The
result factors into two independent theorems: the signal space of a formula is
exactly its set of satisfying signals (Theorem~\ref{thm:A}), and the
canonicalisation of a region of that space depends only on the region and not on
how it was decomposed (Theorem~\ref{thm:B}).
\end{abstract}

\section{What is to be demonstrated}\label{sec:goal}

Write $\Sig(\varphi)$ for the \emph{signal space} of $\varphi$ and $\can$ for
canonicalisation. The goal is
\[
  \varphi \equiv \theta
  \quad\Longleftrightarrow\quad
  \can(\Sig(\varphi)) = \can(\Sig(\theta)),
\]
which factors as
\[
  \varphi \equiv \theta
  \;\underset{\text{Thm.\ \ref{thm:A}}}{\Longleftrightarrow}\;
  \Sig(\varphi) = \Sig(\theta)
  \;\underset{\text{Thm.\ \ref{thm:B}}}{\Longrightarrow}\;
  \can(\Sig(\varphi)) = \can(\Sig(\theta)).
\]

The second step is needed because a signal space arrives \emph{decomposed} into
boxes, and the same region admits many decompositions; two equivalent formulas
generally have different ones. Canonicalisation must therefore be shown to
depend on the region alone.

The first step is a theorem rather than a definition because $\Sig$ is defined
\emph{denotationally}, by structural recursion on the formula
(Section~\ref{sec:eval}). Had $\Sig$ instead been read off a proof search --- a
tableau, say --- Theorem~\ref{thm:A} would additionally require that search to be
sound and complete. Defining $\Sig$ directly reduces it to an induction with one
case per connective.

\section{Preliminaries}

\subsection{Syntax}

Fix a finite set $V$ of variables. Formulas are
\[
  \varphi ::= v \bowtie c \;\mid\; \mathrm{true} \;\mid\; \mathrm{false}
  \;\mid\; \varphi \wedge \varphi \;\mid\; \varphi \vee \varphi
  \;\mid\; F_{[a,b]}\varphi \;\mid\; G_{[a,b]}\varphi
  \;\mid\; \varphi \, U_{[a,b]} \, \varphi
\]
with $v \in V$, $\bowtie \; \in \{<,\le,>,\ge,=,\ne\}$, $a \le b \in \N$, and
$c \in \Q$. Write $\mathrm{vars}(\varphi)$ for the variables occurring in
$\varphi$.

One convention is a choice, not a consequence, and is flagged wherever used:
\textbf{negation is not a constructor}. It is the operation $\lnot(\cdot)$ of
Section~\ref{sec:neg}, mapping a formula to another formula of the same grammar.
Every formula is therefore in negation normal form, and the evaluation of
Section~\ref{sec:eval} has no negation case.

\subsection{Satisfaction}

A \emph{signal} is a map $w : \N \to V \to \R$. Satisfaction $w,t \models
\varphi$ is defined as usual, with
\begin{align*}
  w,t \models v \bowtie c
    &\iff w(t)(v) \bowtie c,\\
  w,t \models F_{[a,b]}\varphi
    &\iff \exists u \in [a,b].\; w,t{+}u \models \varphi,\\
  w,t \models G_{[a,b]}\varphi
    &\iff \forall u \in [a,b].\; w,t{+}u \models \varphi,\\
  w,t \models \varphi \, U_{[a,b]} \, \psi
    &\iff \exists u \in [a,b].\;
       \bigl(w,t{+}u \models \psi \;\wedge\;
             \forall s \in [t,\,t{+}u].\; w,s \models \varphi\bigr),
\end{align*}
and the evident clauses for $\mathrm{true}$, $\mathrm{false}$, $\wedge$, $\vee$.

\begin{remark}[the until convention]\label{rem:until}
The invariant is required on the \emph{closed} $[t, t{+}u]$, hence also at the
witness. Textbook STL uses the half-open $[t, t{+}u)$. The two differ exactly
there: $x{>}0 \; U_{[0,0]} \; y{>}3$ is $x{>}0 \wedge y{>}3$ here and $y{>}3$
in the textbook, and in general $\varphi \, U \, \psi$ here is
$\varphi \, U_{\text{textbook}} \, (\varphi \wedge \psi)$. Both readings are
equally easy to prove correct below --- the case in
Theorem~\ref{thm:A} has the same shape and only the range of the inner
intersection changes.
\end{remark}

\subsection{Horizon and locality}

Let $h(v \bowtie c) = h(\mathrm{true}) = h(\mathrm{false}) = 0$,
$h(\varphi \wedge \psi) = h(\varphi \vee \psi) = \max(h(\varphi),h(\psi))$,
$h(F_{[a,b]}\varphi) = h(G_{[a,b]}\varphi) = b + h(\varphi)$, and
$h(\varphi\,U_{[a,b]}\,\psi) = b + \max(h(\varphi),h(\psi))$.

\begin{lemma}[locality]\label{lem:local}
$w,t \models \varphi$ depends only on $w(s)(v)$ for $s \in [t,\, t + h(\varphi)]$
and $v \in \mathrm{vars}(\varphi)$.
\end{lemma}

\begin{proof}
Induction on $\varphi$, on both indices at once.

\emph{Time.} Atoms and constants use only $s = t$. The Boolean cases
need $[t, t{+}h(\varphi)]$ and $[t, t{+}h(\psi)]$, both inside the maximum. For
$F$ and $G$ the body is evaluated at $t{+}u$ with $u \le b$, needing
$[t{+}u,\, t{+}u{+}h(\varphi)] \subseteq [t,\, t{+}b{+}h(\varphi)]$. For $U$ the
witness needs $[t{+}u,\, t{+}u{+}h(\psi)]$ and the invariant, evaluated at
$s \in [t, t{+}u]$, needs $[s,\, s{+}h(\varphi)]$; both lie in
$[t,\, t{+}h(\varphi\,U\,\psi)]$.

\emph{Variables.} An atom $v \bowtie c$ reads $w(t)(v)$ alone and has
$\mathrm{vars} = \{v\}$; constants read nothing. Every other clause is a Boolean
combination of subformula judgements, and $\mathrm{vars}$ of a compound is the
union of its subformulas', so each hypothesis concerns a subset of
$\mathrm{vars}(\varphi)$.
\end{proof}

Both halves are used below: the time half licenses truncating at $H$, the
variable half licenses an ambient $V$ carrying axes the formula never mentions.

\subsection{The ambient grid, boxes, regions}\label{sec:grid}

Fix $H \in \N$ and finite $V$. The \emph{axis set} is $\Ax = \{0,\dots,H\}
\times V$, and a signal restricted to $\Ax$ is a point of $\R^{\Ax}$.

\begin{quote}
\textbf{Precondition (P1).} $H \ge h(\varphi)$ and $V \supseteq
\mathrm{vars}(\varphi)$.
\end{quote}

Both are needed: with $V \not\supseteq \mathrm{vars}(\varphi)$ the region would
silently lose a variable's constraints and denote a \emph{different} formula,
and with $H < h(\varphi)$ it would be truncated. By Lemma~\ref{lem:local}, if
$H \ge \max(h(\varphi), h(\theta))$ then equivalence over $\R^{\Ax}$ coincides
with equivalence over all signals, so restricting to the grid loses nothing.

An \emph{interval} $I$ is a subset of $\R$ of the form $\{x : x \succ l \text{
and } x \prec r\}$ with $l, r \in \R \cup \{\pm\infty\}$, where each of
$\succ, \prec$ is strict or not independently; infinite endpoints are strict.
$I$ is empty iff $l > r$, or $l = r$ with either end strict.

A \emph{box} is a partial map $\beta$ from $\Ax$ to non-empty intervals,
denoting
\[
  \sem{\beta} = \{ w \in \R^{\Ax} : \forall a \in \dom\beta.\; w(a) \in \beta(a) \};
\]
an axis outside $\dom\beta$ is unconstrained. A \emph{region} is a finite set
$P$ of boxes, denoting $\sem{P} = \bigcup_{\beta \in P} \sem{\beta}$. Thus
$\sem{\emptyset} = \emptyset$ and $\sem{\{\,\emptyset\,\}} = \R^{\Ax}$.

Write $\beta_a$ for $\beta(a)$ when $a \in \dom\beta$ and for $(-\infty,\infty)$
otherwise, so that $\sem{\beta} = \prod_{a \in \Ax} \beta_a$ for every box,
partial or not. Under this convention every box may be treated as total on $\Ax$,
and Section~\ref{sec:B} never needs to assume it.

\begin{lemma}[box algebra]\label{lem:algebra}
Regions are closed under union and intersection. Writing
$(\beta \sqcap \beta')_a = \beta_a \cap \beta'_a$ and
\[
  P \sqcap P' = \{\, \beta \sqcap \beta' \;:\; \beta \in P,\; \beta' \in P',\;
  \beta \sqcap \beta' \neq \emptyset \,\},
\]
one has $\sem{P \cup P'} = \sem{P} \cup \sem{P'}$ and
$\sem{P \sqcap P'} = \sem{P} \cap \sem{P'}$.
\end{lemma}

\begin{proof}
Union is immediate. For a single pair of boxes, membership in both is axis-wise
membership in both intervals, which is membership in the axis-wise intersection;
a product is empty exactly when some factor is. Distributing $\cap$ over $\cup$
gives the general case.
\end{proof}

That intervals carry \emph{independent} endpoint strictness is what makes the
next lemma exact, and it is used throughout: with closed-only intervals the
complement of $[c,\infty)$ would have to be $(-\infty,c]$ and the two would
share $c$.

\section{Theorem A: the signal space is the semantics}

\subsection{Atoms and negation}\label{sec:neg}

\begin{lemma}[atoms]\label{lem:atom}
For $c \in \Q$, the set $\{x \in \R : x \bowtie c\}$ is a union of at most two
intervals: a single interval for $\bowtie \; \in \{<,\le,>,\ge\}$, the point
$[c,c]$ for $=$, and the two strict half-lines $(-\infty,c) \cup (c,\infty)$
for $\ne$. Strict comparisons give strict endpoints, and every endpoint is
either $c$ or infinite.
\end{lemma}

\begin{proof}
By cases; each is the definition of the comparison.
\end{proof}

Let $\overline{\bowtie}$ swap the pairs $(>,\le)$, $(\ge,<)$, $(=,\ne)$.

\begin{lemma}[complementation]\label{lem:neg-op}
$\overline{\;\cdot\;}$ is an involution, and $\{x : x \bowtie c\}$ and
$\{x : x \mathrel{\overline{\bowtie}} c\}$ partition $\R$.
\end{lemma}

\begin{proof}
Each listed pair is swapped, so the operation is an involution; and each pair of
comparisons is complementary on $\R$.
\end{proof}

Define $\lnot(\cdot)$ on formulas by $\lnot(v \bowtie c) = v
\mathrel{\overline{\bowtie}} c$, De Morgan on $\wedge,\vee$, the $F$/$G$
duality, $\lnot\mathrm{true} = \mathrm{false}$ and conversely, and on $U$ by
$\lnot(\varphi\,U\,\psi) = \lnot(\mathrm{exp}(\varphi\,U\,\psi))$ where
$\mathrm{exp}$ is the expansion of Lemma~\ref{lem:expand}.

\begin{lemma}[until is sugar]\label{lem:expand}
$\displaystyle
  w,t \models \varphi\,U_{[a,b]}\,\psi
  \iff
  w,t \models \bigvee_{u=a}^{b}
      \bigl( G_{[0,u]}\varphi \;\wedge\; G_{[u,u]}\psi \bigr).
$
\end{lemma}

\begin{proof}
For fixed $u$: $w,t \models G_{[0,u]}\varphi$ says $w,s \models \varphi$ for all
$s \in [t, t{+}u]$, and $w,t \models G_{[u,u]}\psi$ says $w,t{+}u \models \psi$.
So the $u$-th disjunct holds exactly when $u$ witnesses the until clause;
disjoining over $u \in [a,b]$ matches the existential. (The invariant window is
closed, per Remark~\ref{rem:until}.)
\end{proof}

\begin{lemma}[negation]\label{lem:neg}
$w,t \models \lnot\varphi \iff w,t \not\models \varphi$.
\end{lemma}

\begin{proof}
Let $\mathrm{uk}(\varphi)$ be the maximum nesting depth of $U$ in $\varphi$ and
$|\varphi|$ its size. Induct on $(\mathrm{uk}(\varphi), |\varphi|)$ ordered
lexicographically.

Atoms and constants are Lemma~\ref{lem:neg-op}. The Boolean cases are De
Morgan and the $F$/$G$ cases are the duality of $\exists$ and $\forall$; in each
the hypothesis is applied to strict subformulas, for which $\mathrm{uk}$ does
not increase and $|\cdot|$ strictly decreases.

For $\varphi\,U_{[a,b]}\,\psi$, let $E$ be its expansion. By
Lemma~\ref{lem:expand} it suffices that $w,t \models \lnot E \iff w,t
\not\models E$. Now $E$ is built from $\wedge$, $\vee$, $G$ and copies of
$\varphi$ and $\psi$, so
$\mathrm{uk}(E) = \max(\mathrm{uk}(\varphi), \mathrm{uk}(\psi)) <
\mathrm{uk}(\varphi\,U\,\psi)$ and the hypothesis applies to $E$.
\end{proof}

The measure matters: $\mathrm{exp}$ \emph{increases} size, so an induction on
$|\varphi|$ alone would not be well founded. It strictly decreases $U$-nesting
depth, which is what makes the recursion terminate.

Lemma~\ref{lem:neg} is what licenses the absence of a negation case in
Theorem~\ref{thm:A}. A formula written with $\lnot$ is rewritten by
$\lnot(\cdot)$ before evaluation; the lemma says that rewriting preserves
satisfaction, so the evaluated formula denotes what the written one means. The
saving is not cosmetic: complementing a region of $k$ boxes over $d$ axes costs
$d^{\,k}$ boxes, whereas complementing an atom is one application of
Lemma~\ref{lem:neg-op}.

\begin{lemma}[surface syntax]\label{lem:nnf}
Let the \emph{surface} grammar extend that of Section~\ref{sec:goal} by
$\lnot\varphi$ and $\varphi \to \psi$, and let $\mathrm{nnf}(\cdot)$ map it into
the grammar proper by
\[
  \mathrm{nnf}(\lnot\varphi) = \lnot\,\mathrm{nnf}(\varphi), \qquad
  \mathrm{nnf}(\varphi \to \psi) = \lnot\,\mathrm{nnf}(\varphi) \vee
  \mathrm{nnf}(\psi),
\]
and homomorphically elsewhere. Then $w,t \models \mathrm{nnf}(\varphi) \iff w,t
\models \varphi$, and $\mathrm{nnf}$ preserves $h$ and $\mathrm{vars}$.
\end{lemma}

\begin{proof}
Induction on the surface formula; the homomorphic cases are immediate. For
$\lnot$, the hypothesis gives $w,t \models \mathrm{nnf}(\varphi) \iff w,t \models
\varphi$ and Lemma~\ref{lem:neg} gives $w,t \models \lnot\,\mathrm{nnf}(\varphi)
\iff w,t \not\models \mathrm{nnf}(\varphi)$. For $\to$, $\varphi \to \psi$ holds
iff $\varphi$ fails or $\psi$ holds, which is the displayed disjunction by the
same two facts. Preservation of $h$ and $\mathrm{vars}$ follows from the
hypothesis together with $h(\lnot\varphi) = h(\varphi)$ and
$\mathrm{vars}(\lnot\varphi) = \mathrm{vars}(\varphi)$, which hold since $\lnot$
swaps $\wedge/\vee$ and $F/G$ without changing their bounds, and since both sides
of the $U$ case pass through the expansion $E$ of Lemma~\ref{lem:expand}, for
which $h(E) = b + \max(h(\varphi),h(\psi)) = h(\varphi\,U_{[a,b]}\,\psi)$.
\end{proof}

This is what connects what is written to what is proved: the grammar has neither
$\lnot$ nor $\to$, surface formulas are read up to $\mathrm{nnf}$, and every
result below is a statement about $\mathrm{nnf}(\varphi)$. Preservation of $h$
and $\mathrm{vars}$ is what lets one grid serve a formula and its normal form
alike.

\subsection{Evaluation}\label{sec:eval}

Define a region $\mathcal{E}(\varphi, t)$ by structural recursion:
\begin{align*}
  \mathcal{E}(v \bowtie c,\, t) &= \{\, \{(t,v) \mapsto I\} : I \text{ an interval of Lemma~\ref{lem:atom}} \,\},\\
  \mathcal{E}(\mathrm{true},\, t) &= \{\,\emptyset\,\}, \qquad
  \mathcal{E}(\mathrm{false},\, t) = \emptyset,\\
  \mathcal{E}(\varphi \wedge \psi,\, t) &= \mathcal{E}(\varphi,t) \sqcap \mathcal{E}(\psi,t), \qquad
  \mathcal{E}(\varphi \vee \psi,\, t) = \mathcal{E}(\varphi,t) \cup \mathcal{E}(\psi,t),\\
  \mathcal{E}(F_{[a,b]}\varphi,\, t) &= \bigcup_{u=a}^{b} \mathcal{E}(\varphi,\, t{+}u), \qquad
  \mathcal{E}(G_{[a,b]}\varphi,\, t) = \bigsqcap_{u=a}^{b} \mathcal{E}(\varphi,\, t{+}u),\\
  \mathcal{E}(\varphi\,U_{[a,b]}\,\psi,\, t) &=
     \bigcup_{u=a}^{b}\Bigl( \mathcal{E}(\psi,\, t{+}u) \;\sqcap
       \bigsqcap_{s=t}^{t+u} \mathcal{E}(\varphi,\, s) \Bigr).
\end{align*}

\begin{theorem}\label{thm:A}
Assume \textup{(P1)}. For every $\varphi$ and every $t$ with $t + h(\varphi) \le H$,
\[
  \sem{\mathcal{E}(\varphi, t)} = \{\, w \in \R^{\Ax} : w,t \models \varphi \,\}.
\]
\end{theorem}

\begin{proof}
Structural induction on $\varphi$. There is no negation case: negation is not a
constructor.

\emph{Atoms.} One box per interval of Lemma~\ref{lem:atom}, each constraining
only $(t,v)$; the union of their denotations is $\{w : w(t)(v) \bowtie c\}$.
(P1) guarantees $(t,v) \in \Ax$. Constants are immediate from
$\sem{\{\emptyset\}} = \R^{\Ax}$ and $\sem{\emptyset} = \emptyset$.

\emph{Boolean cases.} Lemma~\ref{lem:algebra} with the induction hypothesis.

\emph{$F$ and $G$.} By Lemma~\ref{lem:algebra} and the hypothesis, the union
over $u \in [a,b]$ denotes $\{w : \exists u.\ w,t{+}u \models \varphi\}$ and the
intersection denotes $\{w : \forall u.\ w,t{+}u \models \varphi\}$; the empty
union is $\emptyset$ and the empty intersection $\R^{\Ax}$, matching the units.
Every instant used satisfies $t{+}u \le t + b \le t + h \le H$ by (P1).

\emph{Until.} For fixed $u$ the term denotes, by Lemma~\ref{lem:algebra} and the
hypothesis,
\[
  \{w : w,t{+}u \models \psi\} \;\cap\; \bigcap_{s \in [t,t+u]} \{w : w,s \models \varphi\},
\]
i.e.\ the signals for which $u$ witnesses the until clause. Taking the union
over $u \in [a,b]$ gives exactly that clause.
\end{proof}

The until case is the \emph{only} place Remark~\ref{rem:until} is used:
replacing the inner index range $[t, t{+}u]$ by $[t, t{+}u)$ proves the textbook
variant instead, with no other change.

\subsection{The signal space and the biconditional}

$\mathcal{E}(\varphi,0)$ is a \emph{decomposition} --- a set of boxes --- and
canonicalisation is defined on decompositions. Its denotation
\[
  \Sig(\varphi) = \sem{\mathcal{E}(\varphi, 0)} \subseteq \R^{\Ax}
\]
is a \emph{region}, a set of signals, and it is regions that
Corollary~\ref{cor:A1} compares. The distinction is kept throughout: $\can$ is
written with a decomposition as argument until Theorem~\ref{thm:B} licenses
otherwise. Write $\varphi \equiv_{\Ax} \theta$ for ``$w,0 \models \varphi$ iff
$w,0 \models \theta$, for all $w \in \R^{\Ax}$''.

\begin{corollary}\label{cor:A1}
Assume \textup{(P1)} for both formulas. Then $\varphi \equiv_{\Ax} \theta$ iff
$\Sig(\varphi) = \Sig(\theta)$.
\end{corollary}

\begin{proof}
By Theorem~\ref{thm:A}, $\Sig(\varphi) = \{w : w,0 \models \varphi\}$ and
likewise for $\theta$. Two subsets of $\R^{\Ax}$ defined by predicates are equal
iff the predicates agree pointwise.
\end{proof}

\begin{corollary}[the grid loses nothing]\label{cor:A2}
Write $\varphi \equiv \theta$ for equivalence at $0$ over \emph{all} signals. If
$H \ge \max(h(\varphi),h(\theta))$ and $V \supseteq \mathrm{vars}(\varphi) \cup
\mathrm{vars}(\theta)$, then $\varphi \equiv \theta$ iff $\varphi \equiv_{\Ax}
\theta$.
\end{corollary}

\begin{proof}
$(\Rightarrow)$ is restriction. $(\Leftarrow)$ Let $w$ be any signal. By
Lemma~\ref{lem:local}, $w,0 \models \varphi$ depends only on $w(s)(v)$ for
$s \le h(\varphi) \le H$ and $v \in \mathrm{vars}(\varphi) \subseteq V$ --- that
is, only on $w|_{\Ax}$; likewise for $\theta$. So $w,0 \models \varphi$ iff
$w|_{\Ax},0 \models \varphi$ iff $w|_{\Ax},0 \models \theta$ iff
$w,0 \models \theta$.
\end{proof}

Both halves of Lemma~\ref{lem:local} are used: the time half to truncate at $H$,
the variable half to allow $V$ to carry axes neither formula mentions. Without
this corollary the development would be about a chosen grid rather than about the
formulas.

\section{Theorem B: canonicalisation depends only on the region}\label{sec:B}

\subsection{Arrangements}

A \emph{decomposition} is a finite set $P$ of boxes over a fixed axis set $\Ax$,
read through the totality convention of Section~\ref{sec:grid}: a box that omits
an axis constrains it by $(-\infty,\infty)$. Nothing below assumes the boxes are
literally total --- $(-\infty,\infty)$ contributes no finite endpoint to $B_a$,
Lemma~\ref{lem:dich} holds vacuously for it, and Lemma~\ref{lem:const} reads
$\beta_a$ uniformly. Write $R = \sem{P}$. For $a \in \Ax$ let $B_a(P)$ be the
finite endpoints on axis $a$ of the intervals of boxes of $P$. Writing
$B_a = \{b_1 < \dots < b_k\}$, the \emph{atoms} of $a$ are
\[
  (-\infty,b_1),\; [b_1,b_1],\; (b_1,b_2),\; [b_2,b_2],\; \dots,\; [b_k,b_k],\; (b_k,\infty),
\]
and $(-\infty,\infty)$ when $B_a = \emptyset$. They are non-empty, pairwise
disjoint and cover $\R$, hence \emph{partition} it. An \emph{atom-product} is
$X = \prod_{a} \alpha_a$ with $\alpha_a$ an atom of $a$; atom-products partition
$\R^{\Ax}$. Write $B(P) = (B_a(P))_{a \in \Ax}$ for the whole family, the
\emph{arrangement} of $P$; the atoms, and hence the atom-products, depend on $P$
only through it.

Including the degenerate atoms $[b_i,b_i]$ is essential. An arrangement using
only $(\dots,b_i)$ and $[b_i,\dots)$ refines a single box but not a union of
boxes that disagree about whether $b_i$ itself belongs: for
$x \ge 1 \wedge y > -1$ and $x > 1 \wedge y \ge -1$, Lemma~\ref{lem:dich} would
fail, since $[1,\infty)$ neither lies inside $(1,\infty)$ nor is disjoint from
it.

\begin{lemma}[dichotomy]\label{lem:dich}
Let $\alpha$ be an atom of axis $a$ and $J$ any interval whose finite endpoints
lie in $B_a$. Then $\alpha \subseteq J$ or $\alpha \cap J = \emptyset$.
\end{lemma}

\begin{proof}
Suppose $x \in \alpha \cap J$. If $\alpha = [b_i,b_i]$ then $\alpha = \{x\}
\subseteq J$. Otherwise $\alpha$ is an open interval $(p,q)$ with $p,q$
consecutive in $B_a \cup \{\pm\infty\}$, so $\alpha$ contains no point of $B_a$.
Let $y \in \alpha$ and suppose $y \notin J$; say $y > x$, the other case being
symmetric. Put $\rho = \sup J$. Then $\rho \le y$, since $x < y < \rho$ would
place $y$ in the interior of $J$; and $\rho \ge x$ since $x \in J$. So $\rho$ is
finite and lies in $[x,y] \subseteq \alpha$, whence $\rho \notin B_a$ ---
contradicting that the endpoints of $J$ lie in $B_a$. (The boundary cases
$\rho = x$ and $\rho = y$ are excluded by the same fact: each would put a point
of $B_a$ inside $\alpha$.) So $\alpha \subseteq J$.
\end{proof}

\begin{definition}[fine cells]\label{def:fine}
$\fine(P) = \{\, X : X \text{ an atom-product},\; X \subseteq R \,\}$.
\end{definition}

The definition mentions $P$ only through the arrangement $B(P)$ and through
$R = \sem{P}$. That two decompositions of one region may nonetheless have
different arrangements, and hence different fine cells, is exactly the difficulty
Theorem~\ref{thm:B} must overcome.

\subsection{Cross-sections}

For $a \in \Ax$ and $x \in \R$ let
\[
  \sigma_a(x) = \{\, y \in \R^{\Ax \setminus \{a\}} : (x,y) \in R \,\}
\]
be the \emph{cross-section} of $R$ at $x$ along $a$. Thus $\alpha \times Z
\subseteq R$ iff $Z \subseteq \sigma_a(x)$ for every $x \in \alpha$.

\begin{lemma}[constancy]\label{lem:const}
$\sigma_a$ is constant on each atom of axis $a$.
\end{lemma}

\begin{proof}
By the totality convention every $\beta \in P$ constrains axis $a$, and
$(x,y) \in R$ iff $y$ lies in the off-$a$ part of some $\beta$ with $x \in
\beta_a$. Let $x, x'$ lie in one atom $\alpha$. By Lemma~\ref{lem:dich},
$x \in \beta_a \iff \alpha \subseteq \beta_a \iff x' \in \beta_a$, so the same
boxes are selected for $x$ and for $x'$, giving the same cross-section.
\end{proof}

This is the exact sense in which a decomposition cannot hide a feature of $R$:
the region is constant across every atom, so wherever $R$ genuinely changes,
\emph{every} decomposition has a breakpoint. It is the only place the
arrangement and the region are related, and everything below rests on it.

\subsection{Runs and the coarse grid}

\begin{definition}[runs]\label{def:run}
For $x \in \R$ with $\sigma_a(x) \neq \emptyset$, the \emph{run} of $x$ on axis
$a$ is
\[
  S_a(x) = \bigcup \{\, I : I \text{ an interval},\; x \in I,\;
  \sigma_a \text{ constant and non-empty on } I \,\}.
\]
\end{definition}

Each $S_a(x)$ is again an interval on which $\sigma_a$ is constant and non-empty:
all the united intervals contain $x$, so their union is an interval, and any of
its points lies in one of them and so has cross-section $\sigma_a(x)$. Hence
$S_a(x)$ is the largest such interval containing $x$. If $S_a(x)$ and $S_a(x')$
meet then their union is an interval on which $\sigma_a$ is constant and
non-empty and which contains both $x$ and $x'$, hence lies in each by
maximality; so they are equal. The runs of axis $a$ are therefore pairwise
disjoint and cover $\{x : \sigma_a(x) \neq \emptyset\}$. They depend on nothing
but $R$.

\begin{corollary}[runs are unions of atoms]\label{cor:runatom}
Every run of axis $a$ is a union of atoms of $a$.
\end{corollary}

\begin{proof}
Let $\alpha$ be an atom meeting a run $S$, say at $z$. By
Lemma~\ref{lem:const}, $\sigma_a$ is constant on $\alpha$ with value
$\sigma_a(z)$, which is the value on $S$ and is non-empty. Since $\alpha$ and $S$
overlap, $\alpha \cup S$ is an interval on which $\sigma_a$ is constant and
non-empty, so $\alpha \cup S \subseteq S_a(z) = S$ by maximality; that is,
$\alpha \subseteq S$.
\end{proof}

The \emph{coarse grid} is the set of products $C = \prod_{a} S_a$ with each
$S_a$ a run of $a$; these are the \emph{coarse boxes}. They are pairwise disjoint,
and by Corollary~\ref{cor:runatom} each is a union of atom-products.

For $X = \prod_a \alpha_a \in \fine(P)$ and each $a$, pick $x \in \alpha_a$ and a
point $z$ of the remaining factors; then $(x,z) \in X \subseteq R$, so
$\sigma_a(x) \neq \emptyset$ and $\alpha_a$ meets the run $S_a(x)$, hence lies
inside it by Corollary~\ref{cor:runatom}. Runs being disjoint, that run is the
unique one containing $\alpha_a$ and so does not depend on the choice of $x$;
writing it $S_a$,
\[
  \mathrm{coarse}(X) = \prod_a S_a
\]
is a well-defined coarse box containing $X$.

\begin{definition}[canonical form]\label{def:canon}
$\can(P) = \{\, \mathrm{coarse}(X) : X \in \fine(P) \,\}$.
\end{definition}

\begin{lemma}[saturation]\label{lem:sat}
If $X \in \fine(P)$ and $X \subseteq C$ for a coarse box $C$, then $C \subseteq R$.
\end{lemma}

\begin{proof}
Write $X = \prod \alpha_a$ and $C = \prod S_a$; as both products are non-empty,
$\alpha_a \subseteq S_a$ for each $a$. Let $Y = \prod \gamma_a$ be any
atom-product with $\gamma_a \subseteq S_a$. Enumerate $\Ax = \{a_1,\dots,a_n\}$,
put $X_0 = X$, and let $X_i$ be $X_{i-1}$ with its $a_i$-component replaced by
$\gamma_{a_i}$. We show $X_i \in \fine(P)$ by induction on $i$.

$X_0 \in \fine(P)$ by assumption. For the step write $a := a_i$, let $\alpha$ be
the $a$-component of $X_{i-1}$ and $\gamma = \gamma_a$, and let $Z$ be the
common product of the other components, so $X_{i-1} = \alpha \times Z$ and
$X_i = \gamma \times Z$. Both $\alpha$ and $\gamma$ lie in the run $S_a$, on
which $\sigma_a$ is constant with some value $\tau$. From $X_{i-1} \subseteq R$
we get $Z \subseteq \sigma_a(x) = \tau$ for $x \in \alpha$; and for $x' \in
\gamma$ we have $\sigma_a(x') = \tau \supseteq Z$, so $\gamma \times Z \subseteq
R$. As $X_i$ is an atom-product, $X_i \in \fine(P)$.

Hence $X_n = Y \in \fine(P)$, so $Y \subseteq R$. As $C$ is the union of such
$Y$ by Corollary~\ref{cor:runatom}, $C \subseteq R$.
\end{proof}

\begin{proposition}\label{prop:charac}
$\can(P) = \{\, C \text{ a coarse box} : C \subseteq R \,\}$.
\end{proposition}

\begin{proof}
$(\subseteq)$ Each element is $\mathrm{coarse}(X)$ for some $X \in \fine(P)$, and
$X \subseteq \mathrm{coarse}(X)$, so Lemma~\ref{lem:sat} applies.
$(\supseteq)$ Let $C$ be a coarse box with $C \subseteq R$. Runs are non-empty,
so $C$ is non-empty and, by Corollary~\ref{cor:runatom}, contains an
atom-product $X$. Then $X \subseteq R$, so $X \in \fine(P)$, and
$\mathrm{coarse}(X) = C$ because on each axis the run containing $X$'s atom is
unique.
\end{proof}

\begin{theorem}\label{thm:B}
Let $P_1, P_2$ be decompositions over the same axis set $\Ax$ with
$\sem{P_1} = \sem{P_2} = R$. Then $\can(P_1) = \can(P_2)$.
\end{theorem}

\begin{proof}
The runs of each axis are defined from $\sigma_a$, hence from $R$ alone, so
$P_1$ and $P_2$ induce the same coarse grid. By Proposition~\ref{prop:charac}
each canonical form is the set of coarse boxes contained in $R$, again a function
of $R$ alone.
\end{proof}

The two decompositions may well have different arrangements, hence different
atoms and different fine cells; what Proposition~\ref{prop:charac} shows is that
this difference cancels. Thus $\can$ factors through the denotation: there is a
well-defined $\can(R)$ with $\can(\sem{P}) = \can(P)$, and every application of
$\can$ to a region rather than a decomposition is licensed by this and by
nothing else.

\begin{proposition}[losslessness]\label{prop:lossless}
$\sem{\can(P)} = R$.
\end{proposition}

\begin{proof}
$(\subseteq)$ Every element of $\can(P)$ lies in $R$ by
Proposition~\ref{prop:charac}.
$(\supseteq)$ Let $p \in R$. Atom-products partition $\R^{\Ax}$, so $p$ lies in
exactly one, say $X$. Choose $\beta \in P$ with $p \in \beta$; by
Lemma~\ref{lem:dich} on each axis, $X$'s atom there meets and hence lies inside
$\beta$'s interval, so $X \subseteq \beta \subseteq R$. Then $X \in \fine(P)$,
and $\mathrm{coarse}(X) \in \can(P)$ contains $X \ni p$.
\end{proof}

\begin{corollary}\label{cor:B1}
Let $\Ax = \{0,\dots,H\} \times V$ with $H \ge \max(h(\varphi),h(\theta))$ and
$V \supseteq \mathrm{vars}(\varphi) \cup \mathrm{vars}(\theta)$, so that
\textup{(P1)} holds for both. Then
\[
  \varphi \equiv \theta
  \iff \varphi \equiv_{\Ax} \theta
  \iff \can(\mathcal{E}(\varphi,0)) = \can(\mathcal{E}(\theta,0)).
\]
\end{corollary}

\begin{proof}
The first equivalence is Corollary~\ref{cor:A2}. For the second:
$(\Rightarrow)$ Corollary~\ref{cor:A1} gives $\Sig(\varphi) = \Sig(\theta)$;
apply Theorem~\ref{thm:B}. $(\Leftarrow)$ By
Proposition~\ref{prop:lossless} each canonical form denotes its own region, so
equal canonical forms give $\Sig(\varphi) = \Sig(\theta)$, and
Corollary~\ref{cor:A1} concludes.
\end{proof}

Canonicalisation is applied to the \emph{decompositions} $\mathcal{E}(\varphi,0)$
and $\mathcal{E}(\theta,0)$, not to the regions; Theorem~\ref{thm:B} is exactly
what permits the shorter $\can(\Sig(\varphi)) = \can(\Sig(\theta))$ of
Section~\ref{sec:goal}, and permits it only after the fact. A single shared
$\Ax$ is required and is not implied by (P1) holding separately for each formula:
the two must be compared on one grid, or $\equiv_{\Ax}$ relates nothing.

For formulas of the surface syntax, Lemma~\ref{lem:nnf} gives
$\varphi \equiv \theta \iff \mathrm{nnf}(\varphi) \equiv \mathrm{nnf}(\theta)$
together with the equal horizons and variable sets that let one $\Ax$ serve both,
so the chain from the formula as written to equal canonical forms is complete.

\section{Remarks}

\begin{remark}[minimality is impossible, not merely absent]
The canonical form is not the smallest cover of $R$ by boxes, and no canonical
construction can be. Take
$R = ([0,2]\times[0,1]) \cup ([0,1]\times[1,2])$. Its canonical form has three
cells, whereas two distinct two-box covers exist:
\[
  [0,2]\times[0,1] \;\cup\; [0,1]\times(1,2]
  \qquad\text{and}\qquad
  [0,1]\times[0,2] \;\cup\; (1,2]\times[0,1].
\]
Both denote $R$, so a minimal cover is not unique and cannot be canonical. What
avoids the choice is that the coarse grid is a \emph{product}: each axis is
partitioned on its own, by a criterion depending only on $R$, and
Lemma~\ref{lem:sat} then quantifies over the whole product at once. Any
construction that instead coalesces cells in pairs must fix an order on the axes,
and the two covers above are what the two orders give.
\end{remark}

\begin{lemma}[provenance of endpoints]\label{lem:prov}
Let $\mathrm{con}(\varphi)$ be the set of constants occurring in the atoms of
$\varphi$. Then every finite endpoint of every interval of every box of
$\mathcal{E}(\varphi,t)$ lies in $\mathrm{con}(\varphi)$; consequently
\[
  \bigcup_{a \in \Ax} B_a(\mathcal{E}(\varphi,0)) \;\subseteq\;
  \mathrm{con}(\varphi) \;\subseteq\; \Q .
\]
\end{lemma}

\begin{proof}
Induction on $\varphi$. For an atom the intervals are those of
Lemma~\ref{lem:atom}, whose finite endpoints are the constant itself; the two
constants contribute no interval at all. Union introduces no interval not already
present in an argument, and $\sqcap$ replaces $\beta_a, \beta'_a$ by $\beta_a
\cap \beta'_a$, whose endpoints are among those of the two (the larger left and
the smaller right). So no clause introduces an endpoint absent from its
arguments.
\end{proof}

Since the atoms of an axis are delimited by $B_a$ and, by
Corollary~\ref{cor:runatom}, runs are unions of atoms, the endpoints of every
coarse box are likewise drawn from $\mathrm{con}(\varphi) \cup \{\pm\infty\}$:
the canonical form of a formula is bounded, syntactically, by the constants the
formula mentions. The lemma is stated for its own interest and is used nowhere
below.

\begin{remark}[hypotheses]
The results above use: \textup{(P1)} $H \ge h(\varphi)$ and $V \supseteq
\mathrm{vars}(\varphi)$; that the two formulas compared share one grid $\Ax$;
that atoms have the form $v \bowtie c$ with $c \in \Q$; that time is discrete and
bounded; and the until convention of Remark~\ref{rem:until}. Totality of boxes is
\emph{not} among them: Section~\ref{sec:grid}'s convention discharges it.
Theorem~\ref{thm:A} assumes nothing beyond the recursion of
Section~\ref{sec:eval}, and Theorem~\ref{thm:B} nothing beyond
Definitions~\ref{def:fine}, \ref{def:run} and \ref{def:canon}.
\end{remark}

\end{document}
